\SourceLeaf{392}%
\[
\begin{aligned}
\mathrm{V.}\quad&2aa\sin2(A''-A')+2a'a'\sin2(A-A'')+2a''a''\sin2(A'-A)\\
&\quad+t[\sin2(A''-A')+\sin2(A-A'')+\sin2(A'-A)]\\
&\quad-4ar\cos(A-\varphi)\sin2(A''-A')\\
&\quad-4a'r\cos(A'-\varphi)\sin2(A-A'')\\
&\quad-4a''r\cos(A''-\varphi)\sin2(A'-A)=0
\end{aligned}
\]
\par
The sum by which \(t\) is multiplied here can also be expressed as
\[
4\sin(A''-A')\sin(A''-A)\sin(A'-A)
\]

Treating equations I, II, IV in a similar way gives an analogous equation VI, which differs from V only in that \(a'',A''\) is replaced by \(a''',A'''\). If we eliminate \(t\) from equations V and VI, we readily see that the result is an equation of the form
\[
\mathrm{VII.}\qquad B+Cr\cos\varphi+Dr\sin\varphi=0
\]
where \(B,C,D\) are known quantities. Their values are easy to work out, but we shall soon show how to avoid doing so. Equation VII makes it clear that the centre of every ellipse tangent to the four sides of our quadrilateral lies on a straight line inclined to the axis of abscissas at an angle with tangent \(=-\frac{C}{D}\), and that the point of intersection has abscissa \(-\frac{B}{C}\). The position of this straight line can, however, be determined much more easily by the following considerations. A diagonal of the quadrilateral can be regarded as a vanishing ellipse tangent to the sides of the quadrilateral. Its centre then clearly lies at the midpoint of the diagonal. It follows readily that the straight line above, which is the locus of the centres of all ellipses tangent to the four sides of the quadrilateral, must be the line joining the midpoints of the two diagonals. It can therefore be found easily. I add two remarks on this point.

1) If the two midpoints coincide (in which case the quadrilateral is a parallelogram), this way of determining the
